\begin{lemma}
Let $G$ be a $\Delta$-regular graph, let $I$ be sampled by
\noderef{RandomIndependentSetSampling} with parameter $\gamma$, fix a vertex
$r$, and let $X\subseteq N_G(r)$.  Assume that no two distinct vertices of
$X$ have a common neighbor outside $N_G[r]$.  Let $\mathcal J_2$ be the
independent pairs in $X$.  For $uv\in\mathcal J_2$, put
$c_{uv}=|N_G(u)\cap N_G(v)|$.  Then
\[
\mathbb E\binom{|I\cap X|}{2}
=
\sum_{uv\in\mathcal J_2}
\frac{2}{\Delta^2}
\int_{0\le x\le y\le\gamma}
\left(1-\frac{x}{\Delta}\right)^{\Delta-c_{uv}}
\left(1-\frac{y}{\Delta}\right)^\Delta\,dy\,dx .
\]
\end{lemma}
\begin{proof}
Let $\overline P_X=\binom{|I\cap X|}{2}$.  Since the output set in
\noderef{RandomIndependentSetSampling} is independent, a two-subset of $X$
contributes to $\overline P_X$ exactly when it is an element of
$\mathcal J_2$ and both of its vertices survive.  Linearity of expectation
therefore gives
\[
\mathbb E\overline P_X
=\sum_{uv\in\mathcal J_2}\Pr[u\in I\text{ and }v\in I].
\]

Fix $uv\in\mathcal J_2$.  Condition on the event that $u$ and $v$ are
activated and use the rescaled priority variables
\[
x=\gamma(1-\pi(u)),\qquad y=\gamma(1-\pi(v)).
\]
The two activation probabilities and the uniform priority law in
\noderef{RandomIndependentSetSampling} give density
$dx\,dy/\Delta^2$ on the square $[0,\gamma]^2$.

Consider first the ordered chamber $x\le y$, equivalently
$\pi(u)\ge\pi(v)$.  A common neighbor of $u$ and $v$ defeats at least one of
the two vertices exactly when it is activated and its priority is larger than
$\pi(v)$ or $\pi(u)$; on this chamber that means priority larger than
$\pi(v)$, and the conditional defeat probability is $y/\Delta$.  Thus each
common neighbor contributes the survival factor $1-y/\Delta$.  A neighbor
of $u$ that is not adjacent to $v$ defeats only $u$, so it contributes
$1-x/\Delta$.  A neighbor of $v$ that is not adjacent to $u$ contributes
$1-y/\Delta$.

Because $G$ is $\Delta$-regular, both $u$ and $v$ have exactly $\Delta$
neighbors.  The clean-neighborhood hypothesis is used here only to ensure
that the common-neighbor count relevant to the two survival events is the
actual graph-theoretic number $c_{uv}=|N_G(u)\cap N_G(v)|$ and that no
outside vertex supplies an additional common-neighbor correlation.  Hence
there are $c_{uv}$ common neighbors, $\Delta-c_{uv}$ neighbors of $u$ not
adjacent to $v$, and $\Delta-c_{uv}$ neighbors of $v$ not adjacent to $u$.
Multiplying the independent survival factors on the chamber $x\le y$ gives
\[
\left(1-\frac{x}{\Delta}\right)^{\Delta-c_{uv}}
\left(1-\frac{y}{\Delta}\right)^{c_{uv}}
\left(1-\frac{y}{\Delta}\right)^{\Delta-c_{uv}}
=
\left(1-\frac{x}{\Delta}\right)^{\Delta-c_{uv}}
\left(1-\frac{y}{\Delta}\right)^\Delta .
\]

On the opposite chamber $y\le x$ the same calculation with $u$ and $v$
interchanged gives
\[
\left(1-\frac{y}{\Delta}\right)^{\Delta-c_{uv}}
\left(1-\frac{x}{\Delta}\right)^\Delta .
\]
The integral over the opposite chamber becomes the same integral after
swapping the names of $x$ and $y$.  Thus the two chambers give two copies
of the integral with $x$ the smaller rescaled priority and $y$ the larger
one:
\[
\Pr[u,v\in I]
=
\frac{2}{\Delta^2}
\int_{0\le x\le y\le\gamma}
\left(1-\frac{x}{\Delta}\right)^{\Delta-c_{uv}}
\left(1-\frac{y}{\Delta}\right)^\Delta\,dy\,dx .
\]
Substitution into the linearity formula proves the asserted identity.
\end{proof}
